\section{Mục tiêu nghiên cứu}
\label{sec:objectives}

Mục tiêu chính của luận văn này là phát triển một phương pháp học sâu hiệu quả để phân loại lối sống thực khuẩn thể từ các contig metagenomic, đặc biệt tập trung vào việc cải thiện hiệu suất trên các đoạn DNA ngắn và trung bình. Cụ thể, luận văn hướng đến các mục tiêu sau:

\subsection{Mục tiêu 1: Phát triển mô hình PhaBERT-CNN}

Thiết kế và triển khai một kiến trúc học sâu lai (hybrid deep learning architecture) kết hợp mô hình nền tảng DNABERT-2 \cite{zhou2023dnabert} đã được tiền huấn luyện với mạng nơ-ron tích chập đa tỷ lệ (multi-scale convolutional neural networks) để phân loại lối sống thực khuẩn thể. Mô hình đề xuất, được đặt tên là PhaBERT-CNN, nhằm tận dụng:

\begin{itemize}
    \item \textbf{Học chuyển giao (Transfer learning):} Sử dụng kiến thức được mã hóa trong DNABERT-2 từ việc tiền huấn luyện trên corpus đa loài (multi-species corpus) để nắm bắt các mẫu tiến hóa ngữ cảnh trong chuỗi DNA của phage
    \item \textbf{Trích xuất đặc trưng đa tỷ lệ:} Sử dụng các nhánh CNN song song với nhiều trường tiếp nhận (receptive fields) để trích xuất các đặc trưng chuỗi cục bộ từ các embedding đã học
    \item \textbf{Tổng hợp thông tin toàn cục:} Áp dụng cơ chế pooling dựa trên attention để tổng hợp thông tin mức chuỗi toàn cục từ các biểu diễn transformer
\end{itemize}

\subsection{Mục tiêu 2: Đánh giá hiệu suất trên nhiều nhóm độ dài contig}

Tiến hành đánh giá toàn diện hiệu suất của PhaBERT-CNN trên bốn nhóm độ dài contig khác nhau:
\begin{itemize}
    \item Nhóm A: 100-400 bp (contig rất ngắn)
    \item Nhóm B: 400-800 bp (contig ngắn)
    \item Nhóm C: 800-1200 bp (contig trung bình)
    \item Nhóm D: 1200-1800 bp (contig dài)
\end{itemize}

Phân tích chi tiết sự thay đổi hiệu suất theo độ dài contig để hiểu rõ điểm mạnh và hạn chế của phương pháp đề xuất trong các kịch bản metagenomic thực tế.

\subsection{Mục tiêu 3: So sánh với các phương pháp tiên tiến}

Thực hiện so sánh có hệ thống giữa PhaBERT-CNN và các phương pháp tiên tiến hiện có, bao gồm:
\begin{itemize}
    \item \textbf{PhaTYP} \cite{shang2023phatyp}: Phương pháp dựa trên BERT với tokenization dựa trên protein
    \item \textbf{DeePhage} \cite{wu2021deephage}: Phương pháp CNN trên chuỗi DNA được mã hóa one-hot
    \item \textbf{ProkBERT} \cite{ligeti2024prokbert}: Phương pháp dựa trên BERT với tokenization k-mer chồng lấp
    \item \textbf{DNABERT-2} \cite{zhou2023dnabert}: Mô hình nền tảng genome với tokenization BPE
\end{itemize}

Phân tích ưu điểm và nhược điểm của từng phương pháp để làm rõ đóng góp của kiến trúc PhaBERT-CNN.

\subsection{Mục tiêu 4: Phân tích tác động của các thành phần kiến trúc}

Nghiên cứu đóng góp của các thành phần khác nhau trong kiến trúc PhaBERT-CNN thông qua:
\begin{itemize}
    \item So sánh PhaBERT-CNN với DNABERT-2 baseline để đánh giá tác động của các lớp CNN và attention pooling
    \item Phân tích sự cân bằng giữa độ nhạy (sensitivity) và độ đặc hiệu (specificity) trên các nhóm độ dài khác nhau
    \item Đánh giá trade-off giữa độ chính xác và thời gian tính toán
\end{itemize}

\subsection{Mục tiêu 5: Đề xuất chiến lược huấn luyện hiệu quả}

Phát triển và đánh giá chiến lược tinh chỉnh hai giai đoạn (two-phase fine-tuning strategy):
\begin{itemize}
    \item \textbf{Giai đoạn warm-up:} Đóng băng backbone DNABERT-2 và chỉ huấn luyện các lớp trích xuất đặc trưng task-specific
    \item \textbf{Giai đoạn tinh chỉnh đầy đủ:} Mở băng toàn bộ kiến trúc và áp dụng learning rate phân biệt (discriminative learning rates) cho các thành phần khác nhau
\end{itemize}

Chiến lược này nhằm đảm bảo sự hội tụ ổn định và thích ứng hiệu quả của mô hình tiền huấn luyện với nhiệm vụ phân loại lối sống phage.

\subsection{Mục tiêu 6: Cung cấp phân tích và hướng dẫn thực tiễn}

Cung cấp phân tích chi tiết về:
\begin{itemize}
    \item Các trường hợp mô hình hoạt động tốt và các trường hợp thất bại
    \item Hạn chế của phương pháp đề xuất và các phương pháp hiện có
    \item Hướng dẫn thực tiễn cho việc lựa chọn phương pháp phù hợp dựa trên đặc điểm của dữ liệu metagenomic
    \item Hướng nghiên cứu tương lai để cải thiện hiệu suất phân loại lối sống phage
\end{itemize}

Thông qua việc đạt được các mục tiêu trên, luận văn kỳ vọng đóng góp vào lĩnh vực sinh tin học (bioinformatics) bằng cách cung cấp một công cụ mạnh mẽ và đáng tin cậy cho phân loại lối sống thực khuẩn thể từ dữ liệu metagenomic, đặc biệt là trên các contig ngắn và trung bình - phạm vi độ dài thách thức nhất trong phân tích metagenomic thực tế.
